\documentclass[11pt]{letter}
\usepackage[utf8]{inputenc}
\usepackage[margin=2.5cm]{geometry}
\usepackage{hyperref}

\address{Rodolfo H. Rodrigo \\
Instituto de Autom\'atica (INAUT) \\
UNSJ--CONICET \\
San Juan J5400ARL, Argentina \\
\href{mailto:rrodrigo@inaut.unsj.edu.ar}{rrodrigo@inaut.unsj.edu.ar}}

\signature{Rodolfo H. Rodrigo \\ on behalf of the authors}

\begin{document}

\begin{letter}{Editor-in-Chief \\ \emph{Applied Mathematics and Computation} \\ Elsevier}

\opening{Dear Editor,}

We are pleased to submit for your consideration the manuscript entitled
\emph{``A Unified Homotopy Framework for Nonlinear System Identification
and Simulation: Constructive Dimensioning of Bilinear Neural
Architectures,''} co-authored by Rodolfo H.\ Rodrigo and H.\ Daniel
Pati\~no (Instituto de Autom\'atica, INAUT, UNSJ--CONICET).

The manuscript develops a unified algorithmic framework for two fundamental
computational problems in nonlinear dynamical systems: parametric function
approximation from sparse input--output data (system identification), and
forward numerical integration of the resulting ordinary differential
equation (simulation). Both tasks are reduced to closed-form Newton ($z_1$)
and Halley ($z_2$) corrections derived from Liao's zeroth-order homotopy
deformation equation, exploiting the bilinear parameter structure shared
by common function representations (radial basis functions, multilayer
perceptrons).

The principal theoretical contribution is a \emph{constructive dimensioning
chain} that maps a target forward-simulation accuracy to approximation-theoretic
bounds on the minimum number of basis functions, the minimum number of
training samples (via Jacobian conditioning), and the maximum admissible
integration time step (via Newton convergence on the discrete residual).
The chain is constructive in the sense that each link is computable from
system quantities, and it admits an operator-theoretic interpretation:
the choice of the auxiliary linear operator $\mathcal{L}$ encodes prior
knowledge that reduces the residual nonlinearity the approximator must
learn, thereby reducing the required approximator dimension.

We validate the framework on three benchmark problems: a one-dimensional
discharge-flow system (orifice-tank model), a two-dimensional Lotka-Volterra
predator-prey system, and robustness experiments under varying levels of
measurement noise. Small-scale parametric approximators (5--8 basis functions)
identified from sparse data (20--40 samples) attain residuals below $10^{-6}$
in a single outer iteration and yield forward-simulation errors below 2\%
over 200-step horizons. We also provide extended wall-clock comparisons
with modern gradient-based optimizers (Adam, L-BFGS) demonstrating the
computational advantages of exploiting bilinear structure in small-network
regimes.

The scope of the empirical validation is deliberately narrow---low-dimensional
ODEs, small-scale approximators---and the role of the present work is to
establish the algorithmic framework, the dimensioning theory, and the
proof-of-concept computational regime. This positioning aligns naturally
with the scope of \emph{Applied Mathematics and Computation}: the
contribution falls at the interface of approximation theory, numerical
analysis for differential equations, and computational methods for system
identification, providing rigorous mathematical foundations with targeted
empirical validation rather than large-scale engineering demonstrations.

We confirm that the manuscript is original, has not been published
elsewhere, and is not under consideration by any other journal. Companion
manuscripts by the authors that share the same conceptual lineage but
address distinct mechanisms (integral Volterra formulation; implicit-model
sample complexity; coupled-system identification) are explicitly delineated
in a dedicated section of the manuscript and are at different stages of
preparation in venues that do not overlap with \emph{Applied Mathematics
and Computation}.

The reference implementations, benchmark scripts, and figure-generation
code are publicly available at
\href{https://github.com/rodrigo1964-ai/unified-homotopy-framework}{github.com/rodrigo1964-ai/unified-homotopy-framework}
and archived at Figshare (DOI:
\href{https://doi.org/10.6084/m9.figshare.31955865}{10.6084/m9.figshare.31955865}).

We thank you in advance for considering our work.

\closing{Yours sincerely,}

\end{letter}

\end{document}
